\documentclass{resume}

\usepackage[left=0.5in,top=0.5in,right=0.5in,bottom=0.5in]{geometry}
\newcommand{\tab}[1]{\hspace{.2667\textwidth}\rlap{#1}} 
\newcommand{\itab}[1]{\hspace{0em}\rlap{#1}}
\name{Thor Lindberg, Cand.it.}
\address{hej@thor.computer \\ \href{https://github.com/thordotcomputer}{github.com/thordotcomputer}}

\begin{document}

%----------------------------------------------------------------------------------------
%	EXECUTIVE SUMMARY
%----------------------------------------------------------------------------------------

\begin{rSection}{Profil Tekst}
    {2+ års erfaring med forretningsanalyse, project management, og softwareudvikling.}
    {En tværfaglig profil med baggrund i design og agil udvikling af IT-systemer og -services, samt et bredt perspektiv på IT og digital transformation for virksomheder, non-profit organisationer og offentlige myndigheder.}
    {Jeg er flydende i både Dansk og Engelsk.}
\end{rSection}

%----------------------------------------------------------------------------------------
%	WORK EXPERIENCE SECTION
%----------------------------------------------------------------------------------------

\begin{rSection}{Erhvervserfaring}
    \href{https://www.nordea.dk}{
        \textbf{IT Project Assistant}, Nordea
    }
    \hfill
    Taastrup, Danmark, 2024
    \begin{itemize}
        \item Ansvarlig for data analyse og administration i access management gruppen.
    \end{itemize}

    \href{https://netcompany.com}{
        \textbf{IT Konsulent}, Netcompany
    }
    \hfill
    København, Danmark, 2022 - 2023
    \begin{itemize}
        %        \item Implementerede server-drevet i18n lokalisation i en React Native og React Native for Web applikation.
        %        \item Designede og implementerede to-vejs XML til Java objekt datamapping.
        \item Medvirkede til softwareimplementering på nye og eksisterende arkitekturer i Java og React Native.
        \item Forfattede specifikationer med UML-diagrammer for projekter i samarbejde med interne teams og eksterne kunder.
        \item Ansvarlig for kontakt til kunder og både design- og tekniske workshops samt Minutes of Meeting (MoM).
              %        \item Arbejdede med GDPR compliance gennem design og udarbejdelse af IT systemer.
        \item \textit{Effekt:} Etablerede kunderelationer i nye områder inden for finans- og regeringssektoren.
        \item \textit{Præstation:} Leverede 2 projekter for 2 kunder, der mødte alle deadlines og leverede udover MVP.
    \end{itemize}

    \href{https://politi.dk/virksomheden/national-enhed-for-saerlig-kriminalitet}{
        \textbf{IT-Administrativ Medhjælper}, National enhed for Særlig Kriminalitet
    }
    \hfill
    København, Danmark, 2021 - 2022
    \begin{itemize}
        \item Styrrede og ydede direkte support til brugere af webportal. Forfattede dokumentation til webportal.
        \item \textit{Effekt:} Forbedrede datakvaliteten ved at hjælpe brugere med at rapportere information nøjagtigt.
    \end{itemize}

    \href{https://itu.dk}{
        \textbf{Teaching Assistant}, IT-Universitetet i København
    }
    \hfill
    København, Danmark, 2021 - 2022
    \begin{itemize}
        \item Assisterede med kurser i UX Design, Eksperimentel Design (JavaScript + MQTT) og Programmering (Python).
    \end{itemize}

    \href{https://www.dominos.dk}{
        \textbf{Shift Leader}, Domino's Pizza
    }
    \hfill
    Lyngby, Danmark, 2021 - 2022
    \begin{itemize}
        \item Primær bager af pizza. Ansvarlig for butiksdrift og administration af lager. Delegerede arbejde til leveringsteam.
    \end{itemize}
\end{rSection}

%----------------------------------------------------------------------------------------
%	EDUCATION SECTION
%----------------------------------------------------------------------------------------

\begin{rSection}{Udannelse}
    \href{https://itu.dk}{
        \textbf{Master of Science i Informationsteknologi}, IT-Universitetet i København
    }
    \hfill
    2020 - 2022
    \begin{itemize}
        %        \item \textit{Thesis:} "Type-Extensible Object Notation: JSON with Syntax for Types" (JavaScript, TDD, UML)
        \item \href{https://learnit.itu.dk/local/coursebase/view.php?ciid=666}{Programming Mobile Applications (React Native, Expo, MapView, Firebase)}
        \item \href{https://learnit.itu.dk/local/coursebase/view.php?ciid=589}{Experimental Design in Practice (Distributed Network Interactions, JavaScript, Event Handlers, MQTT)}
        \item \href{https://learnit.itu.dk/local/coursebase/view.php?ciid=578}{Data Design (Tableau, Data Visualization, Data Graphs, Data-Driven Design)}
        \item \href{https://learnit.itu.dk/local/coursebase/view.php?ciid=487}{Advanced Design Processes (Double Diamond, Design Thinking, User-Centered Design, Prototyping, Arduino)}
        \item \href{https://learnit.itu.dk/local/coursebase/view.php?ciid=624}{UX Design I (Iterative Prototyping, Wireframing, Design Theory, Interface Paradigms, Usability Testing)}
        \item \href{https://learnit.itu.dk/local/coursebase/view.php?ciid=681}{UX Design II (Personas, UX Research, Engagement Design, Critical Reflection)}
        \item \textit{Article:} "Annotation Task Design in Medical Imaging" (LaTeX, Multiple Instance Learning, Multi-Task Learning)
    \end{itemize}

    \href{https://sund.ku.dk}{
        \textbf{Bachelor of Science i Sundhed og Informatik}, Københavns Universitet
    }
    \hfill
    2017 - 2020
    \begin{itemize}
        %       \item \textit{Project:} "Tools for Caretakers of Stroke Victims" (Qualitative Methods, Interviews, Co-Design Workshops)
        \item System Development and Databases (Java, SQL, OOP, Relational Databases, UML, Use Cases, Scrum, XML)
        \item \href{https://kurser.ku.dk/course/sitb17003u/2017-2018}{Statistical Methods (R, Probability Distributions, Estimation, Confidence Intervals, Significance Tests (p-values))}
        \item Mathematical Modelling and Programming (MATLAB, Vectors, Matrices, Linear Algebra, Simulation)
        \item \href{https://kurser.dtu.dk/course/22525}{Medical Image Analysis (MATLAB, Bitmap Image Processing (Matrices), Machine Learning (Blob Analysis))}
        \item \href{https://kurser.ku.dk/course/sbia10210u/2020-2021}{Applied Programming for Biosciences (Python with NumPy, Pandas, Matplotlib)}
        \item IT Security (Encryption/Decryption, Private/Public Keys, Phishing, Physical Security)
        \item Interaction Design (HCI, Affordance, Visibility (Gulf of Execution), Feedback (Gulf of Evaluation))
              %       \item Human Anatomy and Physiology (Cells, Metabolism, Organs, Neurons, Endocrine System, Pathogenesis, Etiology)
    \end{itemize}
\end{rSection}

%----------------------------------------------------------------------------------------
%	CERTIFICATES SECTION
%----------------------------------------------------------------------------------------

\begin{rSection}{Professionelle Certifikater}
    \href{https://www.coursera.org/professional-certificates/meta-front-end-developer}{
        \textbf{Front-End Developer}, Meta
    }
    \hfill
    2024

    \href{https://www.coursera.org/learn/laravel-framework-and-php}{
        \textbf{Mastering Laravel Framework and PHP}, Board Infinity
    }
    \hfill
    2024
\end{rSection}

\end{document}